% Translated from sections/07-conclusion.tex on 2026-09-23. The English file is
% canonical. Change it first, then carry the change here.
\section{Conclusion et artefacts publiés}
\label{sec:conclusion}

\claim{C1}{n'est pas évaluée.} Le jeu de référence existe désormais, avec
\goldPages{}~pages corrigées (section~\ref{sec:protocol}). Mais E1 n'a pas été
exécutée sur ce jeu. L'avantage du modèle affiné sur les modèles publiés reste
donc la prémisse non vérifiée qu'il était au début de ce travail. Ce qui a
changé, c'est la raison. Les données ne manquent plus. Elles restent
seulement inutilisées, et elles penchent.

\claim{C2}{est étayée sur une partie de ses données.} La règle se comporte
comme spécifié sur les lignes où le chemin par décalage a tenu. Sur les
\corrOnFindFallback{} des \corrOffered{}~lignes qui ont emprunté le repli par
recherche, elle comparait la réponse à une ligne que la chaîne avait déjà
corrompue. Ces acceptations ne pèsent donc rien, et l'affirmation repose sur
moins d'un tiers des données recueillies pour elle. Savoir si les
modifications acceptées améliorent le texte est une autre question. Il y faut
une référence pour les pages tirées par E2, que le jeu de référence de C1 ne
couvre pas.

\claim{C3}{est étayée.} L'exactitude au niveau du Livre est de
\locLivreAccuracyWorkbook{} sur la cohorte préenregistrée et de
\locLivreAccuracyCatalogue{} sur la seconde, contre un critère qui demandait
plus de 90\,\%. L'exactitude au niveau du chapitre est nettement plus basse,
à \locChapterAccuracyWorkbook{} et \locChapterAccuracyCatalogue{}. Le
préenregistrement attribuait cet écart au choix de l'étendue, et nommait deux
paramètres à régler avant de conclure. Aucun des deux n'a été réglé. Cette
lecture reste donc une attente confirmée dans son sens.

\subsection{Ce que l'évaluation a changé}

\texttt{MIN\_SCORE} a été vérifié plutôt que choisi. La précision est proche de
100\,\% sur tout le balayage, si bien que la courbe ne désigne pas de seuil. La
figure~\ref{fig:e5-sweep} montre que 0,808 se situe en haut du plateau de
couverture, là où rien n'est perdu. C'est moins que \emph{calibré}, et plus
qu'\emph{hérité}. Les cinq autres constantes du tableau~\ref{tab:thresholds}
n'ont pas été balayées, mais elles sont au moins recensées.

L'évaluation a aussi rectifié le dossier. Quatre chiffres que le projet
avançait en prose ont été recherchés, n'ont pas pu être justifiés, et ont été
retirés. L'un d'eux était un CER du modèle affiné sans jeu de test, sans
journal et sans code derrière lui. Ils sont recensés dans
\file{PROVENANCE.md}.

\subsection{Artefacts publiés}

Ce dépôt contient le code d'évaluation, les enregistrements figés des
exécutions et chaque CSV de résultats. Chaque chiffre se recalcule donc à
partir des entrées versionnées, sans GPU, sans modèle et sans réseau.
\file{data/runs/ocr/} contient les enregistrements d'exécution et de
calibration des \ocrBooks{}~livres, \file{data/runs/correction/} la base de
sondage de E2 et ses verdicts figés, \file{data/runs/matcher/} l'export des
alignements, et \file{data/gold/} le jeu de référence de \goldPages{}~pages,
son manifeste et la convention écrite avant l'annotation. La fiche Hugging
Face de \file{amadis-ft.mlmodel} peut désormais être rédigée, parce que la
comparaison des poids rattache le fichier livré à son point de contrôle
d'entraînement.

Aucune image de page n'est publiée en dehors des planches d'illustration de la
section~\ref{sec:corpus}. Les conditions de réutilisation diffèrent selon
l'institution de conservation et ne sont pas réglées pour l'une d'elles au
moins. \file{data/gold/MANIFEST.csv} consigne le détenteur de chaque page,
afin que chaque cas puisse être réglé séparément. Les transcriptions de
référence sont publiées sous CC-BY-4.0 avec le reste de \file{data/}. Elles
constituent un travail éditorial nouveau et non une reproduction de
l'imprimé, lequel est de toute façon dans le domaine public depuis longtemps.

\subsection{Ce qu'un successeur devrait faire en premier}

Exécuter E1. Le jeu de référence est construit, et plus rien ne s'y oppose.
Y joindre les trois pages transcrites à l'aveugle que demande la
section~\ref{sec:protocol}. Elles coûtent une heure de travail, et elles
transforment les deux faiblesses du jeu en chiffres publiables à côté des taux
d'erreur. Ces faiblesses sont l'absence de plancher de bruit et un biais non
mesuré en faveur du modèle affiné. Corriger ensuite le repli de localisation
des segments, qui coûte aujourd'hui à C2 les deux tiers de ses données.
